\chapter{离散网格与收敛诊断}
\label{ch:grid}

评分核沿厚度用十六条控制体。这一选择进入全部主表，因此必须单独回答：把网格减半或加密，停机进度、1680 s 冷点和热阻方差会移动多少。下面只比较同一初值、同一策略、同一噪声下的离散格式，不引入新的物性。网格选择先于策略比较：八格会把 30 mm 对偶的停机偏差翻号，主表若用八格，交叉表的符号会跟着网格走。

\section{为什么要单独做网格}

显式差分的截断误差随控制体厚度的平方下降，冷点取的是十六条里的最低温度，网格变粗等于把最低温度所在的那一格加厚。加厚之后，同一时刻读到的冷点会偏高，进度会偏早，停机偏差 \(e_{\alpha}\) 会向正方向移动。网格引起的移动要单独报告，交叉表比较的才是策略本身。

S0 在 30 mm、对偶、初温 \(5^\circ\mathrm{C}\)、含水 0.72、均匀场上，把肉糜、土豆、面团各 6 个种子用 16 与 8 个控制体各积一次，得到 18 条配对。另在种子 7、肉糜 30 mm、组合档、对偶、初温 \(8^\circ\mathrm{C}\) 上把控制体数取成 \(n\in\{8,12,16,24,32\}\) 各积一次，用来看 1680 s 这一固定时刻的冷点和进度。后者与演示 CSV 使用同一条件。

\section{S0：十六格相对八格}

表\ref{tab:s0}给出 18 条配对的进度差。差的定义是八格停机进度减去十六格停机进度。中位数 \(0.227\)。肉糜六个种子里有五条停在 \(-0.112\) 附近（十六格）与 \(+0.116\) 附近（八格），差约 \(0.227\)；土豆与面团给出同一量级。停止时刻的移动只有几十秒：十六格约 1716 s，八格约 1748 s。冷点所在格子被加厚之后读数偏高，进度提前饱和到设计线之上。

\begin{table}[htbp]
\centering
\caption{S0：同一初值上 16 格与 8 格的停机进度。\(\Delta e=e_8-e_{16}\)。}
\label{tab:s0}
\begin{tabular}{llrrrr}
\toprule
食品 & 种子 & \(e_{16}\) & \(e_8\) & \(\Delta e\) & \(t_{16}/\mathrm{s}\) \\
\midrule
肉糜 & 0 & \(-0.112\) & \(+0.116\) & 0.227 & 1716 \\
肉糜 & 3 & \(-0.155\) & \(+0.115\) & 0.270 & 1714 \\
土豆 & 0 & \(-0.155\) & \(+0.115\) & 0.270 & 1714 \\
土豆 & 1 & \(-0.111\) & \(+0.116\) & 0.227 & 1718 \\
面团 & 1 & \(-0.155\) & \(+0.115\) & 0.270 & 1714 \\
面团 & 3 & \(-0.111\) & \(+0.116\) & 0.227 & 1718 \\
\bottomrule
\end{tabular}
\end{table}

18 条的中位差与表中列出的典型行一致。图\ref{fig:s0pair}把配对画成十六格对八格的散点，点落在对角线上方：八格一律更正。差堆在 \(0.23\) 与 \(0.27\) 两处，没有靠近零的配对。冷点所在格子被加厚之后读数偏高，进度提前饱和到设计线之上。停止时刻的移动只有几十秒：十六格约 1716 s，八格约 1748 s。符号翻转来自冷点定义，不是策略本身。S0 全部钉在对偶、30 mm、初温 \(5^\circ\mathrm{C}\)、含水 0.72、均匀场，三种食品没有把这个差拉开。规格只说网格减半后决策状态的变化要低于阈值，没有在这个核上给出该阈值的数值，因此结论只能写成：十六条的中位数不能当成网格无关。S1 的句子仍然以十六条为归属。

\begin{figure}[htbp]
\centering
\includegraphics[width=0.78\textwidth]{X_s0_pair.png}
\caption{S0 十八条配对：横轴为十六格停机进度偏差，纵轴为八格。点在对角线上方，粗网格把进度读得更正。}
\label{fig:s0pair}
\end{figure}

\begin{figure}[htbp]
\centering
\includegraphics[width=0.78\textwidth]{X_s0_delta.png}
\caption{S0 进度差 \(\Delta e=e_8-e_{16}\) 的分布。中位数 \(0.227\)。}
\label{fig:s0delta}
\end{figure}

把主对照换成八格，对偶策略在 30 mm 上的中位偏差会从负值翻到正值。策略比较必须固定 \(n=16\)。S0 的 18 条不并进 S1 的 19440 分母。

\section{种子 7 上的五档加密}

表\ref{tab:gridn}是种子 7、肉糜 30 mm、组合档、对偶的五档网格。停机偏差从 \(n=8\) 的 \(+0.150\) 降到 \(n=32\) 的 \(+0.128\)，停止时刻从 1858 s 升到 1868 s，累计加热从 1881 降到 1855。热阻结束方差一律是 0.224：探针是否接受新热阻，与沿厚度分成几格无关。

\begin{table}[htbp]
\centering
\caption{种子 7、肉糜 30 mm、组合档、对偶：控制体数对停机与 1680 s 场的影响。}
\label{tab:gridn}
\begin{tabular}{crrrrr}
\toprule
\(n\) & \(e_{\alpha}\) & \(t_{\mathrm{stop}}/\mathrm{s}\) & \(E\) & \(T_{\mathrm{core}}(1680)\) & \(\alpha_m(1680)\) \\
\midrule
8 & \(+0.150\) & 1858 & 1881 & 73.01 & 0.781 \\
12 & \(+0.149\) & 1862 & 1868 & 71.12 & 0.631 \\
16 & \(+0.145\) & 1866 & 1864 & 70.21 & 0.558 \\
24 & \(+0.136\) & 1866 & 1858 & 69.31 & 0.488 \\
32 & \(+0.128\) & 1868 & 1855 & 68.87 & 0.456 \\
\bottomrule
\end{tabular}
\end{table}

1680 s 这一列的移动大于停机列。八格在该时刻已经把冷点读到 \(73.0^\circ\mathrm{C}\)、进度 0.781；十六格是 \(70.21^\circ\mathrm{C}\)、0.558；三十二格是 \(68.87^\circ\mathrm{C}\)、0.456。粗网格把“快要停机”读得更早。附录\ref{app:trace}与演示 CSV 使用 \(n=16\)，1680 s 仍处于品质加热，请求 0.476，冷点 \(70.21^\circ\mathrm{C}\)。把 73.0 与 0.755 写进十六条的叙述，是把八格轨迹误接到十六格日志上。

图\ref{fig:gridn}把五档的停机偏差、停止时刻、1680 s 冷点和进度画在同一组坐标里。停机偏差随 \(n\) 缓降，1680 s 进度随 \(n\) 下降得更陡。加密从 24 到 32 时，1680 s 冷点只再降 \(0.44\,\mathrm{K}\)，十六格已经落在这一缓坡上。主表固定 \(n=16\) 是工程选择：再加密仍会把该时刻的进度向下推，但策略之间的符号关系保持。

\begin{figure}[htbp]
\centering
\includegraphics[width=0.88\textwidth]{X_grid_n.png}
\caption{种子 7 五档网格。停机偏差缓降；1680 s 进度随加密明显下降。结束方差保持 0.224。}
\label{fig:gridn}
\end{figure}

\section{截断误差与冷点定义}

控制体厚度 \(\Delta x=L/n\)。热扩散显式格式的局部截断与 \(\Delta t\) 和 \(\Delta x^2\) 同阶。时间步固定为 \(2\,\mathrm{s}\)，加密只减小空间步。冷点定义为十六条中的最低温度，加密之后最低温度所在的格子更靠近对称面，读数更低。表\ref{tab:gridn}里 1680 s 表面温度同样随 \(n\) 下降：八格 132.5、十六格 125.7、三十二格 122.9。表面格子变薄，界面通量被更薄的一层热容吸收，瞬态峰值下降。

Fourier 数用半厚 \(L=30\,\mathrm{mm}\) 和 \(\alpha=2.0\times 10^{-7}\,\mathrm{m}^2/\mathrm{s}\) 计算。40 s 探针窗 \(\mathrm{Fo}=8.9\times 10^{-3}\)；1680 s 时 \(\mathrm{Fo}=0.37\)；停机 1866 s 时 \(\mathrm{Fo}=0.41\)。网格误差在 Fourier 数还小于 1 的阶段更显著，因为剖面仍陡，最低温度对格子厚度敏感。保持段结束时剖面已经拉平，五档的终点进度都饱和到 1，网格差被进度上限吃掉。比较必须读停机时刻，终点进度没有分辨率。

\section{与主表的衔接}

S1 与档位研究全部使用十六条。八格会把 30 mm 对偶的停机偏差翻号；把八格的 1680 s 冷点写进十六条的附录，会把“切入保持”的时刻提前约三分钟。时间步固定为 \(2\,\mathrm{s}\)，15 mm 的显式稳定上限约 \(2.20\,\mathrm{s}\)。加密从 24 到 32 时，1680 s 冷点只再降 \(0.44\,\mathrm{K}\)。主表用十六条，是因为八格会翻号，三十二格仍给出同一符号。种子 7 的引用一律以 \(n=16\) 的 CSV 为准：停机 1866 s，\(e_{\alpha}=+0.145\)，1680 s 冷点 \(70.21^\circ\mathrm{C}\)、进度 0.558。演示 CSV 与附录\ref{app:trace} 使用同一条轨迹。
